\documentclass[11pt,letterpaper]{article}
\usepackage[utf8]{inputenc}
\usepackage[T1]{fontenc}
\usepackage[provide=*,spanish]{babel}
\DeclareUnicodeCharacter{201C}{\textquotedblleft}
\DeclareUnicodeCharacter{201D}{\textquotedblright}
\usepackage[margin=2.2cm]{geometry}
\usepackage{amsmath,amssymb,booktabs,array,tabularx}
\usepackage{xcolor,tikz}
\usetikzlibrary{automata,positioning,arrows.meta}
\usepackage{hyperref}
\hypersetup{colorlinks=true,urlcolor=blue}
\usepackage{parskip}
\definecolor{navy}{RGB}{33,54,75}
\newcommand{\datos}[1]{\begin{center}{\Large\bfseries\color{navy}#1}\\[8pt]
Compiladores - Equipo 07\\Yañez Barajas Brandon - Model Designer\\2 de octubre de 2026\end{center}}
\newcommand{\tok}[1]{\texttt{#1}}
\begin{document}
\datos{Gramática libre de contexto y Left Factoring}
\section*{1. Ejemplo seleccionado}
Elegí la entrada del profesor \tok{int a = 10;}. También consideré \tok{int a;} para mostrar dos alternativas con un prefijo común. La CFG describe la secuencia de tokens: $id$ representa un identificador y $num$ una constante numérica. Ambos son terminales de esta gramática.

\section*{2. Gramática original}
Defino $G_0=(N_0,T,P_0,S)$, donde:
\[N_0=\{S\},\qquad T=\{\texttt{int},id,num,\texttt{=},\texttt{;}\}.\]
El símbolo inicial es $S$ y las producciones son:
\[S\to\texttt{int}\ id\ \texttt{=}\ num\ \texttt{;}\ \mid\ \texttt{int}\ id\ \texttt{;}.\]
Cada producción tiene un único no terminal a la izquierda; por ello, la gramática es libre de contexto.

\section*{3. Factorización por la izquierda}
Las alternativas comparten el prefijo $\texttt{int}\ id$. Separé este prefijo e introduje el no terminal $R$ para representar el resto:
\[\begin{aligned}S&\to\texttt{int}\ id\ R,\\R&\to\texttt{=}\ num\ \texttt{;}\ \mid\ \texttt{;}.\end{aligned}\]
La nueva gramática es $G_1=(N_1,T,P_1,S)$, con $N_1=\{S,R\}$ y $P_1$ formado por las reglas anteriores. $T$ y el símbolo inicial se conservan. $R\to\texttt{;}$ no es una producción vacía: genera el punto y coma de la declaración sin inicialización.

\section*{4. Derivaciones}
Con inicialización:
\[S\Rightarrow\texttt{int}\ id\ R\Rightarrow\texttt{int}\ id\ \texttt{=}\ num\ \texttt{;}.\]
Sin inicialización:
\[S\Rightarrow\texttt{int}\ id\ R\Rightarrow\texttt{int}\ id\ \texttt{;}.\]
El lexer establece que \tok{a} corresponde a $id$ y \tok{10} a $num$. No se agregan reglas $id\to a$ ni $num\to 10$: la gramática opera sobre categorías de tokens y terminales específicos.

\newpage
\section*{5. Árbol de derivación}
\begin{center}\begin{tikzpicture}[level distance=1.5cm,level 1/.style={sibling distance=3cm},level 2/.style={sibling distance=1.5cm},every node/.style={text=navy}]
\node {$S$} child {node {\texttt{int}}} child {node {$id$}} child {node {$R$} child {node {\texttt{=}}} child {node {$num$}} child {node {\texttt{;}}}};
\end{tikzpicture}\end{center}
Las hojas de izquierda a derecha son $\texttt{int},id,\texttt{=},num,\texttt{;}$. Coinciden con la secuencia de la declaración seleccionada.

\section*{6. Comparación y conservación del lenguaje}
La gramática original genera las dos estructuras siguientes:
\[L(G_0)=\{\texttt{int}\ id\ \texttt{=}\ num\ \texttt{;},\quad\texttt{int}\ id\ \texttt{;}\}.\]
En la gramática factorizada, $S$ genera el prefijo común y $R$ permite exactamente los dos sufijos originales. Por lo tanto, $L(G_0)=L(G_1)$. La transformación reorganiza las producciones sin añadir ni eliminar estructuras.

\section*{7. Relación con el proyecto}
El lexer se implementará mediante expresiones regulares. La CFG describe el orden de los tokens y no depende del mecanismo utilizado para reconocerlos. Los AFND y AFD se conservan como documentación del modelo teórico.

Esta propuesta cubre declaraciones sencillas de tipo \tok{int}; no constituye una gramática completa de C. La tarea presenta la gramática, su transformación y comprobación, sin implicar que el lexer tenga un parser o valide la sintaxis de las entradas.

\end{document}
